\begin{thebibliography}{32}

\bibitem{gottesman2019guidelines}
Gottesman, O., Johansson, F., Komorowski, M., Faisal, A., Sontag, D., Doshi-Velez, F., \& Celi, L. A. (2019). Guidelines for reinforcement learning in healthcare. \textit{Nature Medicine}, 25(1), 16-18.

\bibitem{kumar2020cql}
Kumar, A., Zhou, A., Tucker, G., \& Levine, S. (2020). Conservative q-learning for offline reinforcement learning. \textit{Advances in Neural Information Processing Systems}, 33, 1179-1191.

\bibitem{kostrikov2021iql}
Kostrikov, I., Nair, A., \& Levine, S. (2021). Offline reinforcement learning with implicit q-learning. \textit{arXiv preprint arXiv:2110.06169}.

\bibitem{futoma2020myth}
Futoma, J., Simons, M., Panch, T., Doshi-Velez, F., \& Celi, L. A. (2020). The myth of generalisability in clinical risk prediction: algorithms prepare for the past and are surprised by the future. \textit{The Lancet Digital Health}, 2(9), e489-e492.

\bibitem{trella2024online}
Trella, A. L., et al. (2024). Online reinforcement learning with dynamic fidelity: Balancing exploration and clinical safety in digital interventions. \textit{Journal of Machine Learning Research}, 25(42), 1-34.

\bibitem{parbhoo2022causal}
Parbhoo, S., et al. (2022). Causal off-policy evaluation with latent confounders in clinical sequences. \textit{Advances in Neural Information Processing Systems}, 35.

\bibitem{hang2025suicide}
Hang, J., et al. (2025). Forecasting suicide risk from sparse digital phenotyping data: A multi-task Gaussian Process approach. \textit{Lancet Psychiatry}, 12(2), 112-121.

\bibitem{lu2026neural}
Lu, C., et al. (2026). Neural stochastic differential equations for continuous-time suicide risk forecasting. \textit{Nature Machine Intelligence}, 8(1), 45-56.

\bibitem{nock2026smartphone}
Nock, M. K., et al. (2026). Real-time smartphone monitoring of suicidal thoughts and behaviors: A large-scale clinical validation. \textit{JAMA Psychiatry}, 83(4), 312-320.

\bibitem{cai2026causal}
Cai, X., \& Onnela, J. P. (2026). Causal estimands for n-of-1 digital interventions with informative missingness. \textit{Biometrika}, 113(1), 89-104.

\bibitem{huang2025adolescent}
Huang, J., et al. (2025). Continuous monitoring of adolescent bipolar disorder using passive smartphone sensing. \textit{Neuropsychopharmacology}, 50(3), 441-450.

\bibitem{tu2024amie}
Tu, T., et al. (2024). Towards conversational diagnostic AI. \textit{arXiv preprint arXiv:2401.05654} (DeepMind).

\bibitem{saab2024medgemini}
Saab, K., et al. (2024). Capabilities of Gemini Models in Medicine. \textit{arXiv preprint arXiv:2404.18416}.

\bibitem{chen2018neuralode}
Chen, R. T., Rubanova, Y., Bettencourt, J., \& Duvenaud, D. K. (2018). Neural ordinary differential equations. \textit{Advances in Neural Information Processing Systems}, 31.

\bibitem{debrouwer2019gruode}
De Brouwer, E., Simm, J., Arany, A., \& Moreau, Y. (2019). GRU-ODE-Bayes: Continuous-time RNN with Bayesian update. \textit{Advances in Neural Information Processing Systems}, 32.

\bibitem{lipton2016missing}
Lipton, Z. C., Kale, D., \& Wetzel, R. (2016). Directly modeling missing data in health care with recurrent neural networks. \textit{arXiv preprint arXiv:1606.06758}.

\bibitem{insel2017digital}
Insel, T. R. (2017). Digital phenotyping: technology for a new science of behavior. \textit{World Psychiatry}, 16(2), 170-171.

\bibitem{torous2016digital}
Torous, J., et al. (2016). New tools for new research in psychiatry: a focused review of smartphones and mobile health. \textit{Current Psychiatry Reports}, 18(10), 91.

\bibitem{shalit2017causal}
Shalit, U., Johansson, F. D., \& Sontag, D. (2017). Estimating individual treatment effects using deep latent variable models. \textit{International Conference on Machine Learning (ICML)}, 3076-3085.

\bibitem{kidger2021neuralsde}
Kidger, J., Foster, J., Li, X., \& Lyons, T. (2021). Neural SDEs as infinite-dimensional GANs. \textit{International Conference on Machine Learning (ICML)}, 5454-5463.

\bibitem{fischer2025clinician}
Fischer, A. S., et al. (2025). Aligning RL agents with clinician expectations: A study on antidepressant titration. \textit{Nature Medicine}, 31(1), 58-67.

\bibitem{joshi2021defer}
Joshi, S., et al. (2021). Learning to defer to clinicians in psychiatric emergency departments. \textit{AAAI Conference on Artificial Intelligence}, 35(1), 452-460.

\bibitem{mozannar2025causal}
Mozannar, H., \& Sontag, D. (2025). Causal falsification of reinforcement learning policies in healthcare. \textit{Conference on Health, Inference, and Learning (CHIL)}.

\bibitem{mortimer2025retfound}
Mortimer, B., et al. (2025). RETFound-Green: Energy-efficient foundation models for longitudinal psychiatric sensing. \textit{The Lancet Digital Health}, 7(3), e142-e151.

\bibitem{levine2020offline}
Levine, S., Kumar, A., Tucker, G., \& Fu, J. (2020). Offline reinforcement learning: Tutorial, review, and perspectives on open problems. \textit{arXiv preprint arXiv:2005.01643}.

\bibitem{lin2018}
Lin, H. W., Mermelstein, R. J., \& Hedeker, D. (2018). Latent trait shared-parameter mixed models for missing ecological momentary assessment data. \textit{Statistical Methods in Medical Research}, 27(6), 1845-1857.

\bibitem{li2020}
Li, X., Wong, T.-K. L., Chen, R. T. Q., \& Duvenaud, D. (2020). Scalable Gradients for Stochastic Differential Equations. \textit{Proceedings of the 23rd International Conference on Artificial Intelligence and Statistics (AISTATS)}, PMLR 108:3870-3882.

\bibitem{wu2019}
Wu, Y., Tucker, G., \& Nachum, O. (2019). Behavior Regularized Actor Critic (BRAC). \textit{arXiv preprint arXiv:1911.11361}.

\bibitem{kidger2025compact}
Kidger, P., et al. (2025). Neural Stochastic Differential Equations on Compact State-Spaces. \textit{arXiv preprint arXiv:2508.23000}.

\bibitem{dro2024}
Author, A., et al. (2024). A Unified Distributionally Robust Formulation for Offline RL. \textit{arXiv preprint}.

\end{thebibliography}
